\chapter{Discussion}
\label{chap:discussion}

Chapter \ref{chapter:evaluation} examined whether the proposed LGVI-SDP framework 
preserves the geometric structure
of the modeled rigid-body systems,
produces certifiable open-loop trajectories,
and can generate feasible trajectories for non-tight relaxations in a closed-loop implementation.
The results provide a qualified answer to the research question in Section \ref{sec:introduction_problem_statement}.
The framework is suitable for Lagrangian variational integrators of single- and multibody systems
and can certify particular discrete trajectory-optimization instances.
Still tightness, feasible extraction,
and computational cost depend strongly on the considered dynamics and remain challenging.
In the following, the main findings are discussed and related to the state of the art.
Finally, potential applications are elaborated.

\section{Discussion of Main Findings}
\label{sec:discussion_of_main_findings}
The simulation results in Section \ref{sec:numerical_simulation_results} 
establish that the derived discrete models of the 
3D pendulum on $\mathrm{SO}(3)$ and the underactuated Acrobot $(\mathrm{SO}(2) \times \mathbb{R}^2)^2$
represent the intended dynamics while exhibiting energy-conservation properties over a long time horizon.
For the unforced 3D pendulum,
Figure \ref{fig:simulation_results_3d_pendulum} shows that the LGVI remains on $\mathrm{SO}(3)$,
preserves the momentum
and exhibits bounded long-term energy error.
In comparison, a standard RK4 integrator has errors in all three quantities,
whereas projected RK4 restores the rotation constraint but accelerates energy drift.
The constrained Acrobot results in Figures \ref{fig:simulation_results_acrobot_unforced} 
and \ref{fig:simulation_results_acrobot_forced} extend this observation to a multibody system,
where both link rotations remain on $\mathrm{SO}(2)$,
the holonomic link constraints are satisfied,
and the unforced energy error remains bounded.
These results demonstrate that the derived discrete Lagrangian variational 
integrators of the 3D-pendulum \eqref{eq:forced_3D_pendulum_EL} and the 
Acrobot \eqref{eq:acrobot_full_del_shifted} preserve the relevant configuration 
and constraint geometry over the simulated horizon.


\paragraph{Single-body optimization}
The 3D pendulum demonstrates the complete method for a single rigid body.
Its LGVI can be written as the quadratic POP in \eqref{eq:controlled_3d_pendulum_pop},
but a dense second-order relaxation would already require a $91{,}378\times91{,}378$ 
moment matrix for $N=20$.
The minimum-degree chordal extension generated by SPOT in Section \ref{subsec:sparse_relaxations} 
reduces the largest matrix to $325\times325$,
compared with $820\times820$ for the complete temporal clique.
Therefore, the Moment-SOS relaxation can convert a nonconvex
algebraic LGVI formulation into a sparse convex optimization problem.
Still, the MD sparsity construction is required for computational application.

For all five targets $\theta_{x,des}=\{0^\circ, 45^\circ, 90^\circ, 135^\circ, 180^\circ\}$,
Table \ref{tab:pendulum_3d_sdp_statistics} reports a rank metric below $7.3\times10^{-8}$,
a manifold violation below $4.6\times10^{-8}$,
and a suboptimality gap below $1.8\times10^{-6}$.
Additionally, the solution extractions remain feasible.
Therefore, these values indicate that the second-order relaxations are numerically tight and
certify global optimality for the considered discrete POP.
Figure \ref{fig:optimization_results_x180_md} illustrates the $\theta_{x,des}=180^\circ$ case,
where the extracted trajectory reaches the desired attitude while remaining on $\mathrm{SO}(3)$.
Consequently, the single-body problem of a 3D pendulum on $SO(3)$ is globally certifiable 
using discrete LGVI dynamics and sparse SDP relaxations.

Still, it has to be noted that the certified solutions do not necessarily 
reach the desired states after exactly $N$ steps, while still being certifiable. For example,
the $\theta_{x,des}=90^\circ$ case reaches a terminal state of $\theta_{x,final}\approx70.5^\circ$,
since the implemented target is not a hard constraint but a cost inside the objective function.
Therefore, the certified solution is optimal for the discrete objective, 
but not for reaching the exact terminal state.

\paragraph{Multibody formulation and sparsity}
The Acrobot extends the framework to a constrained and underactuated multibody system.
Its maximal-coordinate LGVI represents both links separately
and imposes the joints through holonomic constraints and multipliers.
This enables a quadratic polynomial structure.
Moreover, the reduction introduced in Section \ref{sec:discrete_motion_planning_acrobot} eliminates 
the position variables and expresses the velocities through rotation variables.
Thereby, the size of the decision vector, which has to be optimized, 
decreases from $17N+3$ to $13N-1$ variables while preserving quadratic polynomial dynamics.
It therefore reduces the SDP without introducing densely coupled dynamical 
representations when using minimal coordinates \cite{BrudigamManchester2021LQR}.

Nevertheless, even after the reduction,
the dense second-order moment matrix has the dimension $304{,}590\times304{,}590$ for $N=60$.
The complete chain-like sparsity construction \eqref{eq:acrobot_full_self_initial_clique}-
\eqref{eq:acrobot_full_self_interior_cliques} replaces the dense 
matrix with separate blocks of sizes $351\times351$ and $253\times253$.
The multibody-specific construction in \eqref{eq:acrobot_hybrid_self_cliques} 
keeps the complete cliques at $k=1,2$,
but separates later dynamics and kinematics into $171\times171$ and $91\times91$ matrices.
The results of the SDP relaxation, shown in Tables \ref{tab:acrobot_sdp_statistics_full} and
\ref{tab:acrobot_sdp_statistics_hybrid},
identify similar relaxation metrics for both structures.
Still, the specialized cliques reduce the MOSEK runtime by a factor of approximately three to five.
Consequently, the identification of problem-specific sparsity is necessary 
for solving multi-body systems using semidefinite programming to obtain 
fast and numerically reliable results \cite{Kang_2025Global}. 

However, it is important to note that multibody-specific sparsity does not guarantee tighter relaxations 
than chain-like sparsity. Therefore, before applying the proposed framework to new multibody systems, 
it is necessary to evaluate the tightness of the relaxation for the specific problem.

\paragraph{Open-loop Acrobot optimization}
The open-loop SDP evaluation in Section \ref{subsec:acrobot_optimization_results} 
separates easy from difficult Acrobot problems.
For the targets $\theta_{\mathrm A}^{\mathrm{des}}=0^\circ$ and
$\theta_{\mathrm A}^{\mathrm{des}}=45^\circ$,
Tables \ref{tab:acrobot_sdp_statistics_full} and \ref{tab:acrobot_sdp_statistics_hybrid} report 
rank-one moment matrices, feasible extractions, and negligible suboptimality gaps.
Therefore, in these cases, the solutions are certified.
However, similar to the 3D-Pendulum, the certified $45^\circ$ solution retains a terminal 
error of approximately $55^\circ$, because the target is only weighted in the objective.

For the targets $\theta_{\mathrm A}^{\mathrm{des}}=90^\circ$,
$\theta_{\mathrm A}^{\mathrm{des}}=135^\circ$, and
$\theta_{\mathrm A}^{\mathrm{des}}=180^\circ$,
the maximum rank metric increases to approximately $0.35$-$0.66$,
while the extracted rotations have manifold violations around $1.0$-$1.4$.
Therefore, the first moments do not define feasible $\mathrm{SO}(2)$ trajectories
and no certificate is obtained. The results indicate that the additional complexity of the multibody
underactuated system increases the difficulty of the convex optimization problem.

For the runtime, the optimization using specialized multibody cliques \eqref{eq:acrobot_hybrid_self_cliques} 
requires approximately $15$-$31$ minutes,
compared with $45$-$113$ minutes for the chain-like cliques \eqref{eq:acrobot_full_self_initial_clique}-
\eqref{eq:acrobot_full_self_interior_cliques}.
Therefore, this demonstrates that the size of the moment matrices is the driving factor for the runtime of the SDP solver.
Still, both optimization variants remain too slow for online control.
A lower relaxation order would reduce matrix sizes but weaken the tightness of the relaxation.
Moreover, instead of using the second order interior-point solver MOSEK \cite{MOSEK}, 
a first-order SDP solver, for example CuADMM \cite{Kang_2026Fast}, could accelerate the same relaxation
at the possible cost of less accurate certified solutions. Additionally, in this thesis, only 
second order $\kappa = 2$ relaxations are considered, as \cite{Sangli2025}
states that second order relaxations are sufficient for most of the considered problems. In general, 
higher order relaxations can be used to tighten the relaxation \eqref{eq:convergence_of_moment_sos_hierarchy}, 
but at the cost of increased computational complexity.

\paragraph{Closed-loop Acrobot stabilization}
Tables \ref{tab:acrobot_sdp_statistics_full} and \ref{tab:acrobot_sdp_statistics_hybrid} 
show that $\delta_1$ is smaller than the worst-horizon metric $\delta_{\max}$ inside 
the open-loop optimization.
This motivates the receding-horizon implementation in Section 
\ref{sec:acrobot_trajectory_optimization_SDP_MPC}, where
only the first control from \eqref{eq:acrobot_mpc_applied_control} is applied.
The remaining non-tight prediction is discarded,
and the SDP is solved again from the simulated state.
However, $\delta_1$ is not a rank-one condition and therefore not tight.
Still, the implemented closed-loop control has a feedback structure, 
which can stabilize the system even if the open-loop solution is not tight 
(see Section \ref{subsec:acrobot_sdp_mpc_results}).

Figure \ref{fig:acrobot_mpc_u5_closed_loop} 
illustrates that this procedure stabilizes the system starting from an initial rotation of $5^\circ$,
after approximately $11.9\,\mathrm{s}$ with $u_{\max}=5\,\mathrm{N\,m}$.
The simulated trajectory remains on $\mathrm{SO}(2)$
although the predicted blocks are not rank one throughout the whole maneuver.
This demonstrates numerical closed-loop feasibility for the tested configuration.
However, the solution cannot be considered globally optimal, since the 
optimization problem is not tight. Additionally, failures are also 
consistent with limited control torque availability for $u_{\max}=5\,\mathrm{N\,m}$.

With $u_{\max}=15\,\mathrm{N\,m}$,
Figure \ref{fig:acrobot_basin_comparison} 
shows that the matched discretization ($\Delta t_{SDP} = \Delta t_{\mathrm{sim}}= \Delta t_{ctrl}$) 
Cases~I and~II stabilize $94$ of $100$ initial configurations.
Therefore, consistently refining all $\Delta t_{SDP}$ and $\Delta t_{\mathrm{sim}}$ intervals from $0.1\,\mathrm{s}$ to $0.05\,\mathrm{s}$
does not change the success rate. Nevertheless,
the different failed trials for the same initial conditions between Cases~I and~II
indicate that the success is not guaranteed for all $\Delta t_{SDP} = \Delta t_{\mathrm{sim}}= \Delta t_{ctrl}$
configurations.

In Case~III, where there is a discretization mismatch 
between the SDP and simulation time steps
($\Delta t_{SDP} = 0.1\,\mathrm{s}$, 
 $\Delta t_{\mathrm{sim}}= 0.0025\,\mathrm{s}$, $\Delta t_{ctrl} = 0.025\,\mathrm{s}$),
the controller stabilizes only $34$ of $100$ configurations.
Consequently, Figure \ref{fig:acrobot_sampled_trajectory_comparison}(b) shows a more oscillatory behaviour
near the target configuration before stabilization, compared to
Case~I in Figure \ref{fig:acrobot_sampled_trajectory_comparison}(a) for the same initial condition of $(\theta_1^{\mathrm{init}},\theta_2^{\mathrm{init}})=(20^\circ,140^\circ)$.
Here, Case~III requires $4.1\,\mathrm{s}$
instead of $2.2\,\mathrm{s}$ in Case~I.
The results indicate that mismatched discretization can lead to 
reduced reliability for stabilization of the Acrobot.

For the successful $(\theta_1^{\mathrm{init}},\theta_2^{\mathrm{init}})=(80^\circ,80^\circ)$ example,
Cases~I, II, and III require approximately $3.1\,\mathrm{h}$,
$27.3\,\mathrm{h}$,
and $17.8\,\mathrm{h}$.
These illustrative runtimes depend on $\Delta t_{\mathrm{SDP}}$ and $N_{\mathrm{SDP}}$,
which determine each SDP size,
and on the number of replanning iterations.
A matched discretization case with $\Delta t_{\mathrm{SDP}} =\Delta t_{\mathrm{sim}}=\Delta t_{\mathrm{ctrl}}=0.0025\,\mathrm{s}$ 
would require $N_{\mathrm{SDP}}=800$, which is not computationally feasible with current hardware.
Consequently, the closed-loop implementation is only considered as an offline application 
instead of a real-time controller. However, the 
open-loop optimization using semidefinite relaxation is also not real-time capable, since the SDP solver requires
approximately $15$-$31$ minutes for the considered $N_{\mathrm{SDP}}=60$ problem.

\section{Comparison with the State of the Art}
\label{sec:comparison_state_of_the_art}

The first contribution of this thesis is the Acrobot-specific LGVI in maximal coordinates
\eqref{eq:acrobot_full_del_shifted}, including its underactuation
\eqref{eq:forced_acrobot_rotational_equations} and holonomic joint constraints. This formulation aligns with
the preservation properties of
\cite{Marsden2001,Lee2008} and the maximal-coordinate formulations in
\cite{BrudigamManchester2021VariationalIntegrators,BrudigamManchester2021LQR}.
Moreover, the reduction in Section \ref{sec:discrete_motion_planning_acrobot}
resulted in a smaller decision vector $\bm{z}$ and a
quadratic polynomial model suitable for SDP relaxation \eqref{eq:controlled_acrobot_pop}.

The closest framework to this thesis is \cite{Sangli2025}. In that work, all reported 
trajectories are passed through a local-search step using IPOPT, including cases for which 
the SDP relaxation provides a small optimality gap or satisfies the rank condition. 
In contrast, the second-order relaxations of all considered 3D-pendulum problems in this 
thesis are numerically rank one and directly yield feasible trajectories without subsequent 
NLP refinement. The agreement between the SDP lower bound and the objective of the extracted 
trajectory therefore certifies global optimality using the SDP solution alone. This 
demonstrates that, for the investigated single-body benchmark, the complete pipeline from the 
nonconvex LGVI formulation to a feasible and globally optimal trajectory can be performed 
without depending on the initialization or convergence of a local solver, in contrast to \cite{Sangli2025}. 
Nevertheless, the drone problems in 
\cite{Sangli2025} include translational dynamics and obstacle avoidance and are therefore 
more complex than the present pendulum problem. Whether similarly direct SDP extraction 
remains possible after introducing obstacle constraints for the 3D-pendulum 
has to be evaluated in future research.
%%%%%%%%%%%%%%%
With respect to the research gap identified in Section~\ref{sec:introduction_research_gap}, this thesis 
demonstrates that the complete LGVI-SDP framework can be transferred from a single rigid body 
to a constrained and underactuated multibody system of increased dynamical complexity. 
The results confirm that globally optimal trajectories can be 
certified when the relaxation is tight and that closed-loop feedback controls can still be 
extracted for some non-tight problems, for example the Acrobot stabilization. 
However, the transfer is not equally reliable for 
all configurations, since relaxation tightness, solver feasibility, and computational cost 
remain strongly problem-dependent. In particular, the closed-loop trajectories do not 
provide certificates of global optimality, and their computation is 
too slow for real-time application. The research gap is therefore addressed for the 
investigated single- and multibody models, but remains open for larger systems, contact 
and obstacle constraints, uncertain dynamics, and hardware experiments.


\section{Potential Applications}
\label{sec:potential_applications}

The most direct application is offline trajectory generation. Tight SDP solutions
can provide certified benchmarks for evaluating local optimizers, as initial guesses 
for related NLP problems, 
or as reference trajectories for online tracking controllers.

For non-tight problems, the SDP-MPC implementation can still generate feasible
realized trajectories. The resulting state-control sequences could
be stored as demonstrations for faster online MPC methods or used to construct
warm starts from previously solved configurations. This is consistent with
\cite{Kang_2026Fast}, where data-based warm starts substantially reduce repeated
SDP solution times.

The trajectories could also provide training data for imitation learning or
offline reinforcement learning. This is particularly relevant for underactuated
systems such as the Acrobot, where swing-up requires coordinated actions over a
longer time horizon and can be difficult to explore \cite{acrobot_RL}.
A learned controller could approximate the SDP closed-loop policy while operating
at lower online cost.
